% Imporditud: tools/import_eksperimendid.py
% Allikas: Eesti füüsikaolümpiaadi eksperimendi ülesannete
%   kogu (Rael Kalda, 2023), ülesanne Ü99, lahendus L99.
\setAuthor{EFO žürii}
\setRound{piirkonnavoor}
\setYear{2009}
\setNumber{G E1}
\setDifficulty{3}
\setTopic{vedelike mehaanika}
\setVahendid{kumminiit, tundmatust materjalist mutter, statiiv, joonlaud, tops veega}
\setPunktid{10}
\setAllikas{Eksperimendi kogu Ü99, lahendus lk 77}
\setLahendusseis{toores}

\prob{Kumminiit}
Leida mutri tihedus. Mõõtevea hindamist ei nõuta.

\hint

\solu
Elastsusjõud Fe = k $\cdot \Delta$l, kus $\Delta$l = (l $-$ $l_0$) on kumminiidi pikenemine (1 p). (Anda see punkt ka siis, kui avaldis ei ole ilmutatult välja kirjutatud, kuid hiljem õigesti kasutatud). Kaaluda me mutrit küll saame, aga arvulist tulemust see meile ei anna, vaid ainult suhtelise pikenemise. Seepärast valemist $\rho$ = m/V lähtuda ei anna. Archimedese seaduse kohaselt väheneb keha kaal vees aga sama palju, kui keha tõrjub vett välja ehk arvuliselt $\rho$vesi $\cdot$ V , kui keha on veest tihedam. Kaalumine õhus annab:

$\rho$mutter $\cdot$ Vmutter $\cdot$ g = k $\cdot$ (l $-$ $l_0$) , (1 p)

kus $l_0$ on kumminiidilõigu algpikkus ja l on kumminiidi pikkus koormatuna mutriga. Kaalumine vees omakorda annab:

($\rho$mutter $- \rho$vesi) $\cdot$ Vmutter $\cdot$ g = k $\cdot$ (lv $-$ $l_0$) , (1,5 p)

kus lv on kumminiidilõigu pikkus mutri kaalumisel vette uputatuna. Jagades võrrandid omavahel läbi, saame

l $-$ $l_0$ = lv $-$ $l_0$ ,

$\rho$mutter $\rho$mutter $- \rho$vesi

kust l $-$ $l_0$ l $-$ lvesi $\rho$mutter = $\cdot \rho$vesi. (1,5 p)

Kõik vajalikud mõõtmised on korralikult teostatud ja dokumenteeritud --- 2 p. Kum-

miniidi pikkus on valitud nii, et l ulatub pea joonlaua maksimummõõduni --- 1 p. Lõppvastus vahemikus 1,90$-$2,20 g/cm$^3$ (2 p) või vahemikus kuni 1,80$-$2,30 g/cm$^3$ (1 p). (Tallinn: lõppvastus vahemikus 2,15 $-$ 2,35 g/cm$^3$ (2 p) või vahemikus kuni 2,05 $-$ 2,45 g/cm$^3$ (1 p).)

Mutri ruumala mõõtmise/hindamise eest punkte üldjuhul mitte anda.

Kui õpilane oma lahenduses vaid oletab, et materjal on alumiinium, ning selle põhjal paneb vastuseks $\rho$ = 2,7 g/cm$^3$, anda kogu lahenduse eest 1 p.

Vale elastsusjõu valemi kasutamine Fe = k$\cdot$l, kus l on niidi pikkus venitatud olekus --- max 5 p.

Märkus: Konkreetne katse on konkreetse kummi ja selle kummi kindla lõigu puhul hästi korratav (mm täpsusega) ja seega korduvad mõõtmised midagi juurde ei anna.

Märkus 2: Mutri materjali tegelik tihedus on 2,7 g/cm$^3$, kuid kuna kasutatav mudel ei võta arvesse kõiki faktoreid, on vastus nihkes.
\probend
